\chapter{Results and Discussion}
\label{chap:results}

This chapter reports the complete, unvarnished empirical evidence of Capstone Project \textbf{FA26AI16} across four evaluation pillars corresponding directly to Review 3 Stages 1, 3, and 4: (1) SOTA baseline reproduction and five-seed forensic auditing (\texttt{B1} and \texttt{B2}); (2) controlled seven-branch ablations (\texttt{A0--A6}) across \texttt{YOLO11s} and \texttt{YOLOv8s} on both validation (316 images) and official test (1,587 images) splits; (3) Before-vs-After inference optimization on RTX 3050---including our dedicated \nolinkurl{RepConv.switch_to_deploy()} equivalence benchmark---and multi-stream edge hardware profiling; and (4) software verification (108 \texttt{pytest} cases), FastAPI load/stress testing, field CCTV testing, and the UAT/SUS protocol. Validation and official test splits are kept strictly separated throughout.

\section{Results and Analysis}

\subsection{SOTA Baselines Reproduction and Forensic Audit (B1 \& B2)}

\subsubsection{Comparison of Published Papers vs. GDUT-HWD Five-Class Reproductions}

Table~\ref{tab:ch6_b1_b2_table1} presents the formal cross-protocol overview contrasting the original SHWD two-class metrics published by Zhao et al.~\cite{zhao2025yolodcrcf} (\texttt{B1}) and Song et al.~\cite{song2024improved} (\texttt{B2}) against Thanh's clean GDUT-HWD five-class reproductions (\texttt{seed = 3407}, Stage 1 fine-tune, \texttt{imgsz = 640}, \texttt{batch = 16}, \texttt{AdamW}, \texttt{lr0 = 1e-4}, dataset fingerprint \texttt{026a2afb...}). Because the published papers evaluated two-class SHWD (with \texttt{B1} trained from scratch for 300 SGD epochs), direct arithmetic subtraction against five-class GDUT scores is invalid; all architectural comparisons are strictly conducted within the identical five-class GDUT split.

\begin{table}[htbp]
\centering
\caption{Table 1 (Baseline Audit): Original Published SOTA Metrics (SHWD 2-Class) vs. Clean GDUT-HWD 5-Class Reproductions (Thanh, \texttt{seed = 3407}).}
\label{tab:ch6_b1_b2_table1}
\resizebox{\textwidth}{!}{%
\renewcommand{\arraystretch}{1.22}
\begin{tabular}{|l|l|r|r|r|r|r|p{3.8cm}|}
\hline
\rowcolor{slatehdr}
\textbf{Model} & \textbf{Source / Split} & \textbf{Params (M)} & \textbf{P (\%)} & \textbf{R (\%)} & \textbf{mAP@50 (\%)} & \textbf{mAP@50--95 (\%)} & \textbf{Protocol / Class \texttt{none}} \\
\hline
\textbf{B1: YOLO-DCRCF} & Paper Zhao 2025 (SHWD 2-cls) & -- (13.4G) & 90.90 & 88.50 & 92.70 & \textbf{60.60} & Scratch 300 ep, SGD \\ \hline
\textbf{B1: YOLO11n-DCRCF} (Thanh) & GDUT Val (316) / \textbf{Test (1,587)} & 6.802 & 89.85 / \textbf{89.57} & 81.53 / \textbf{76.96} & 88.86 / \textbf{84.36} & 57.47 / \textbf{54.37} & Test \texttt{none}: AP50--95 \textbf{47.97\%}, R \textbf{69.90\%} \\ \hline
\textbf{B2: YOLOv8s-DWR-ASPP-NWD} & Paper Song 2024 (SHWD 2-cls) & -- & 91.80 & 86.60 & 92.00 & \textbf{62.20} & 100 ep, SGD, batch 16 \\ \hline
\textbf{B2: YOLOv8s-DWR-ASPP-NWD} (Thanh) & GDUT Val (316) / \textbf{Test (1,587)} & 16.285 & 90.35 / \textbf{90.43} & 86.90 / \textbf{81.81} & 92.13 / \textbf{88.19} & 61.06 / \textbf{57.98} & Test \texttt{none}: AP50--95 \textbf{52.62\%}, R \textbf{76.80\%} \\ \hline
\end{tabular}
}
\end{table}

As an auxiliary sanity check matching Zhao et al.'s original scratch regime on five-class GDUT, training \texttt{B1} from scratch for 300 SGD epochs (best epoch 138) yielded only \textbf{52.02\%} test mAP@50--95 (P $88.74\%$, R $74.85\%$), confirming that Stage 1 warm-start ($54.37\%$ test) is strictly superior to scratch training on 1,265--2,530 images.

\subsubsection{Five-Seed Matched Stability Audit and B1 Tuning Investigation (Dung)}

All 18 core baseline modules (\nolinkurl{baseline_model.py} \texttt{e2dd3140}, \nolinkurl{baseline_modules.py} \texttt{cf8d4639}, \texttt{warm\_model.py} \texttt{a31f59c7}, \texttt{gdut\_data.py} \texttt{620c14fb}) matched SHA-256 byte-for-byte between Thanh and Dung. Table~\ref{tab:ch6_b1_b2_table2} compares Thanh's seed-3407 reproductions against Dung's best \texttt{DCRCF} run (\texttt{Pair-3407 Control}) and five-seed \texttt{V8F} sweep on GDUT Validation (316 images), while Table~\ref{tab:ch6_dcrcf_trials} summarizes every controlled tuning intervention attempted on \texttt{B1} alongside matched five-seed stock controls (\texttt{C0}, \texttt{V8S}).

\begin{table}[htbp]
\centering
\caption{Table 2 (Baseline Audit): Thanh vs. Dung Controlled Validation Comparison (GDUT Val, 316 Images).}
\label{tab:ch6_b1_b2_table2}
\resizebox{\textwidth}{!}{%
\renewcommand{\arraystretch}{1.22}
\begin{tabular}{|l|l|c|r|r|r|r|p{3.9cm}|}
\hline
\rowcolor{slatehdr}
\textbf{Role} & \textbf{Model / Source} & \textbf{Seed} & \textbf{P (\%)} & \textbf{R (\%)} & \textbf{mAP@50 (\%)} & \textbf{mAP@50--95 (\%)} & \textbf{Class \texttt{none} (Val)} \\
\hline
\textbf{B1 (SOTA)} & \texttt{YOLO11n-DCRCF} (Thanh) & 3407 & 89.85 & \textbf{81.53} & 88.86 & 57.47 & AP50--95: \textbf{51.40\%}, R@0.25: \textbf{75.99\%} \\ \hline
\textbf{B1 (Dung Best)} & \texttt{YOLO11n-DCRCF} (\texttt{Pair Control}) & 3407 & \textbf{91.65} & 78.77 & 88.86 & \textbf{57.78} & AP50--95: 51.08\%, R@0.287: 75.74\% (P 83.61\%) \\ \hline
\textbf{B2 (SOTA)} & \texttt{YOLOv8s-DWR-ASPP-NWD} (Thanh) & 3407 & \textbf{90.35} & \textbf{86.90} & \textbf{92.13} & \textbf{61.06} & Single-seed anchor (Test: 57.98\%) \\ \hline
\textbf{B2 (Dung 5-Seed)} & \texttt{YOLOv8s-DWR-ASPP-NWD} (\texttt{V8F}) & 3407--3411 & 90.25 $\pm$ 0.94 & 86.47 $\pm$ 0.58 & 91.82 $\pm$ 0.17 & 60.78 $\pm$ 0.35 & AP50--95: \textbf{55.01\%}, R@0.25: \textbf{83.93\%} (P 84.18\%) \\ \hline
\end{tabular}
}
\end{table}

\begin{table}[htbp]
\centering
\caption{Complete Controlled Tuning Trials on \texttt{B1 (YOLO11n-DCRCF)} and Five-Seed Matched Stock Audits.}
\label{tab:ch6_dcrcf_trials}
\resizebox{\textwidth}{!}{%
\renewcommand{\arraystretch}{1.20}
\begin{tabular}{|p{3.0cm}|p{3.3cm}|p{5.0cm}|p{3.2cm}|}
\hline
\rowcolor{slatehdr}
\textbf{Trial / Run ID} & \textbf{Controlled Intervention} & \textbf{Measured Validation Outcome (316 Images)} & \textbf{Audit Verdict} \\
\hline
\textbf{Trial A} (Eval-only replay) & Re-score Thanh \texttt{B1} ckpt; sweep \texttt{conf} $0.25 \to 0.20$ & Exact match $57.4731\%$ ($<3\times 10^{-6}\%$ delta); lower conf adds $+10$ TP but $+15$ FP. & Threshold tuning cannot fix P3 drop. \\ \hline
\textbf{Trial B} (Split LR pilot) & Two param groups: $1\times 10^{-4}$ old, $3\times 10^{-4}$ new \texttt{SBA/DCNv2} & Overall mAP@50--95 $+0.13$ pp, but \texttt{none} Recall drops \textbf{$-2.97$ pp} ($306 \to 294$ TP). & \textbf{FAIL} -- harms bare-head safety recall. \\ \hline
\textbf{Canonical \texttt{V5}} (\texttt{tuned}) & Global \texttt{lr0 = 3e-4} across full network & Val $57.75\%$ ($+0.28$ pp vs. Thanh); \texttt{none} R@0.25 rises to $77.72\%$ ($+8$ TP), but \texttt{none} FP jumps $61 \to 88$ ($+27$ FP, P $78.11\%$) and total R falls $-1.18$ pp ($81.53\% \to 80.35\%$). & Precision--Recall trade-off only; no net gain. \\ \hline
\textbf{\texttt{Pair-3407}} (\texttt{Control} vs. \texttt{Treatment}) & \nolinkurl{GradScaler(init_scale=16.0)}; \texttt{Treatment} adds \texttt{pos\_weight=1.5} on \texttt{none} & \texttt{Control}: $57.78\%$ (P $+1.80$ pp to $91.65\%$, R $-2.76$ pp to $78.77\%$); \texttt{Treatment}: drops to $57.28\%$ ($-0.50$ pp; \texttt{none} AP $51.08\% \to 50.41\%$). & \texttt{pos\_weight} fails $\le 0.10$ pp gate; \texttt{Control} lowers safety Recall. \\ \hline
\textbf{\texttt{C3} vs. \texttt{C0}} (5 seeds \texttt{3407--3411}) & Matched 5-seed \texttt{YOLO11n-DCRCF} vs. Stock \texttt{YOLO11n} & \texttt{C3}: $57.283 \pm 0.173\%$ vs. Stock \texttt{C0}: \textbf{$58.812 \pm 0.082\%$} (\textbf{$-1.529$ pp} loss; \texttt{none} AP $51.16\%$ vs. $52.93\%$). & \textbf{Definitive:} DCRCF underperforms Stock on GDUT. \\ \hline
\textbf{\texttt{V8F} vs. \texttt{V8S}} (5 seeds \texttt{3407--3411}) & Matched 5-seed \texttt{YOLOv8s-DWR-ASPP-NWD} vs. Stock \texttt{YOLOv8s} & \texttt{V8F}: $60.781 \pm 0.353\%$ vs. Stock \texttt{V8S}: \textbf{$61.414 \pm 0.074\%$} (\textbf{$-0.633$ pp} loss, losing 4/5 seeds; mAP@50 loses 5/5 seeds $-0.488$ pp). & \textbf{Preserved unmodified} as SOTA Baseline 2. \\ \hline
\end{tabular}
}
\end{table}

\subsubsection{Forensic Decomposition of 98 Bare-Head False Negatives in B1 and Root Causes}

Replaying Thanh's \texttt{B1} checkpoint over all 404 ground-truth \texttt{none} boxes across the 316 validation images at \texttt{conf = 0.25} yielded 306 True Positives, 61 False Positives, and \textbf{98 False Negatives} (Figure~\ref{fig:ch6_b1_b2_audit}). Forensically decomposing all 98 misses pinpoints the structural failure mode:
\begin{itemize}[leftmargin=1.8em, itemsep=2pt]
  \item \textbf{84 / 98 misses (85.7\%):} The detector already localized a valid candidate box ($\mathrm{IoU} \ge 0.50$), yet its \texttt{none} classification score collapsed below $0.25$ (68 cases where all five classes fell below $0.25$; 16 cases where one of the four helmet colors outscored \texttt{none}).
  \item \textbf{11 / 98 misses (11.2\%):} Zero predicted candidates reached $\mathrm{IoU} \ge 0.50$---and \textbf{100\% of these 11 misses were tiny bare heads ($< 32$ px)}. Across all 404 \texttt{none} boxes, small boxes ($< 32$ px) accounted for 68/230 misses versus only 30/174 for larger heads.
  \item \textbf{3 / 98 misses (3.1\%):} Valid above-threshold candidates suppressed during NMS.
  \item \textbf{P3 Feature Bottleneck:} \textbf{358 / 393} localized \texttt{none} targets---and \textbf{77 / 87} candidate-bearing false negatives---originated at the high-resolution \textbf{P3 scale ($80 \times 80$, \texttt{model.14})}.
\end{itemize}

Four root causes explain why hyperparameter tuning on \texttt{DCRCF} cannot rescue its ceiling on five-class GDUT: (1) \textbf{Representation discontinuity}---replacing the entire neck with four randomly initialized \texttt{SBA} modules ($4.08\text{M}$ params, $60.03\%$ of network) leaves only $26.22\%$ ($1.783\text{M}$) of Stage 1 weights inherited while forcing the inherited \texttt{Detect 19} head to consume low-amplitude \texttt{SBA} outputs ($\text{RMS} = 0.079\text{--}0.107$); (2) \textbf{Protocol mismatch}---Zhao 2025 trained 300 SGD epochs from scratch on 2 classes, whereas 5-class AdamW fine-tuning peaks early at epoch 33 before validation box loss rises ($1.225 \to 1.273$); (3) \textbf{Non-deterministic \texttt{DCNv2} gradients}---166 CUDA \nolinkurl{compute_grad_input} warnings; and (4) \textbf{Ceiling deficit vs. Stock}---\texttt{C3} ($6.80\text{M}$, $14.8\text{G}$ FLOPs) trails lightweight Stock \texttt{YOLO11n} ($2.59\text{M}$, $6.3\text{G}$) by $-1.53$ pp.

Similarly, for \textbf{Baseline 2 (\texttt{B2})}, preserving $65.16\%$ of Stage 1 weights avoids \texttt{B1}'s feature shock and achieves strong \texttt{none} recall ($83.93\%$ Val on \texttt{V8F}, $76.80\%$ Test on Thanh) thanks to \texttt{NWD}'s small-box center sensitivity. However, the fixed $0.5 \cdot \mathrm{NWD}(C=12.8)$ penalty loosens high-IoU ($0.75\text{--}0.95$) box tightness while \texttt{DWR + ASPP} inflates parameters to \textbf{16.285M ($+46.2\%$)}, compute to \textbf{41.23 GFLOPs ($+45.0\%$)}, and inference latency to \textbf{16.44 ms}. We therefore retain \texttt{B2} untouched as our second faithful SOTA reference. Figures~\ref{fig:ch6_b1_b2_audit} and~\ref{fig:ch6_b1_authentic} illustrate the five-seed audit, 98-FN decomposition, and authentic Kaggle PR/confusion plots from \texttt{gdut-yolo11n-dcrcf-baseline}.

\begin{figure}[htbp]
\centering
\includegraphics[width=0.98\textwidth]{figures/fig_gdut_baselines_b1_b2_audit.png}
\caption{Forensic SOTA baseline audit: (a) Five-seed matched validation stability (\texttt{seeds 3407--3411}) proving both \texttt{B1} (\texttt{C3}) and \texttt{B2} (\texttt{V8F}) trail matched stock YOLO models (\texttt{C0}, \texttt{V8S}); (b) Forensic decomposition of all 98 bare-head (\texttt{none}) misses in \texttt{B1}.}
\label{fig:ch6_b1_b2_audit}
\end{figure}

\begin{figure}[htbp]
\centering
\begin{subfigure}[b]{0.49\textwidth}
  \centering
  \includegraphics[width=\textwidth]{figures/b1_dung_val_pr.png}
  \caption{B1 (\texttt{Pair-3407 Control}) Validation PR Curve}
\end{subfigure}
\hfill
\begin{subfigure}[b]{0.49\textwidth}
  \centering
  \includegraphics[width=\textwidth]{figures/b1_dung_val_cm.png}
  \caption{B1 (\texttt{Pair-3407 Control}) Normalized Confusion Matrix}
\end{subfigure}
\caption{Authentic Kaggle validation plots from Dung's best \texttt{YOLO11n-DCRCF} run (\nolinkurl{tundng111/gdut-yolo11n-dcrcf-baseline}, \texttt{v372678468}), highlighting elevated background false-negative leakage on class \texttt{none}.}
\label{fig:ch6_b1_authentic}
\end{figure}

\subsection{Comprehensive Ablation Study (A0--A6 on YOLO11s \& YOLOv8s)}

All fourteen checkpoints (\texttt{A0/S1} and \texttt{A1--A6} across \texttt{YOLO11s} and \texttt{YOLOv8s}) share the identical cryptographic dataset fingerprint \nolinkurl{ce2a883e0b5050edb96d494218e85a9d16ea04b9a19dc5723c572cb502edfe9c}, deterministic \texttt{seed = 3407}, and identical backbone-specific Stage 1 initialization checkpoints. Validation numbers reflect exact JSON values from \nolinkurl{run_manifest.json} (\texttt{val\_check}, 316 images, 1,738 boxes) rounded to two decimals; Official Test numbers reflect the exact three-decimal Ultralytics logs from \nolinkurl{GDUT_final_test_evaluation.ipynb} (1,587 images, 9,574 boxes, \texttt{split=test}) expressed to one decimal percentage point.

\subsubsection{Validation Split Performance (316 Images, 1,738 Boxes)}

Tables~\ref{tab:ch6_val_y11}, \ref{tab:ch6_val_v8}, and~\ref{tab:ch6_val_deltas} report complete validation metrics and arithmetic percentage-point (pp) deltas across both backbones.

\begin{table}[htbp]
\centering
\caption{GDUT-HWD Validation Results (316 Images) for \texttt{YOLO11s} Across Ablations \texttt{A0--A6}.}
\label{tab:ch6_val_y11}
\small
\renewcommand{\arraystretch}{1.18}
\resizebox{\textwidth}{!}{%
\begin{tabular}{|c|l|r|r|r|r|}
\hline
\rowcolor{slatehdr}
\textbf{Config} & \textbf{Active Modifications (From Stage 1)} & \textbf{Precision (\%)} & \textbf{Recall (\%)} & \textbf{mAP@50 (\%)} & \textbf{mAP@50--95 (\%)} \\
\hline
\textbf{A0} & Stage 1 Stock \texttt{YOLO11s} (base train only) & \textbf{92.21} & 84.80 & 91.83 & 60.83 \\ \hline
\textbf{A1} & Stock + HardAug + \nolinkurl{mixup=0.1, hsv_v=0.45} & 91.33 & 87.37 & \textbf{93.40} & \textbf{61.69} \\ \hline
\textbf{A2} & + \texttt{CoordConv} (\texttt{layer 0}) & 91.01 & 86.25 & 92.72 & 61.31 \\ \hline
\textbf{A3} & + 7$\times$\texttt{RepConv} (\texttt{0,1,3,5,7,17,20}) & 90.79 & \textbf{87.86} & 92.88 & 61.41 \\ \hline
\textbf{A4} & + \texttt{Focal-EIoU} ($\gamma = 0.5$) & 90.64 & 87.49 & 92.77 & 61.36 \\ \hline
\textbf{A5} & + \texttt{BiFormerBlockLite} (\texttt{layer 10}, repl. \texttt{C2PSA}) & 90.05 & 87.14 & 92.40 & 61.08 \\ \hline
\rowcolor{slatehdr}
\textbf{A6} & \textbf{Full Architecture (Coord + 6$\times$Rep + BiFormer + F-EIoU)} & 90.45 & 87.58 & 92.50 & 61.23 \\ \hline
\end{tabular}}
\end{table}

\begin{table}[htbp]
\centering
\caption{GDUT-HWD Validation Results (316 Images) for \texttt{YOLOv8s} Across Ablations \texttt{A0--A6}.}
\label{tab:ch6_val_v8}
\small
\renewcommand{\arraystretch}{1.18}
\resizebox{\textwidth}{!}{%
\begin{tabular}{|c|l|r|r|r|r|}
\hline
\rowcolor{slatehdr}
\textbf{Config} & \textbf{Active Modifications (From Stage 1)} & \textbf{Precision (\%)} & \textbf{Recall (\%)} & \textbf{mAP@50 (\%)} & \textbf{mAP@50--95 (\%)} \\
\hline
\textbf{A0} & Stage 1 Stock \texttt{YOLOv8s} (base train only) & 91.54 & 85.75 & 91.88 & 60.55 \\ \hline
\textbf{A1} & Stock + HardAug + \nolinkurl{mixup=0.1, hsv_v=0.45} & 91.55 & 86.21 & \textbf{93.14} & \textbf{61.19} \\ \hline
\textbf{A2} & + \texttt{CoordConv} (\texttt{layer 0}) & 91.51 & \textbf{86.79} & 92.47 & 60.32 \\ \hline
\textbf{A3} & + 7$\times$\texttt{RepConv} (\texttt{0,1,3,5,7,16,19}) & 91.47 & 86.67 & 92.41 & 60.30 \\ \hline
\textbf{A4} & + \texttt{Focal-EIoU} ($\gamma = 0.5$) & 92.60 & 84.91 & 92.35 & 60.21 \\ \hline
\textbf{A5} & + \texttt{BiFormerBlockLite} (\texttt{layer 9}, after \texttt{SPPF}) & 92.02 & 85.56 & 92.59 & 60.26 \\ \hline
\rowcolor{slatehdr}
\textbf{A6} & \textbf{Full Architecture (Coord + 6$\times$Rep + BiFormer + F-EIoU)} & \textbf{93.21} & 85.11 & 92.56 & 60.68 \\ \hline
\end{tabular}}
\end{table}

\begin{table}[htbp]
\centering
\caption{Validation Percentage-Point (pp) Deltas of Full Architecture \texttt{A6} Against \texttt{A0} and \texttt{A1}.}
\label{tab:ch6_val_deltas}
\small
\renewcommand{\arraystretch}{1.16}
\begin{tabular}{|l|l|r|r|}
\hline
\rowcolor{slatehdr}
\textbf{Comparison} & \textbf{Backbone} & \textbf{$\Delta$ mAP@50 (pp)} & \textbf{$\Delta$ mAP@50--95 (pp)} \\
\hline
\texttt{A6 $-$ A0} & \texttt{YOLO11s} & $+0.67$ & $+0.39$ \\ \hline
\texttt{A6 $-$ A0} & \texttt{YOLOv8s} & $+0.68$ & $+0.13$ \\ \hline
\texttt{A6 $-$ A1} & \texttt{YOLO11s} & $-0.90$ & $-0.46$ \\ \hline
\texttt{A6 $-$ A1} & \texttt{YOLOv8s} & $-0.58$ & $-0.51$ \\ \hline
\end{tabular}
\end{table}

\subsubsection{Official Test Split Performance (1,587 Images, 9,574 Boxes)}

Tables~\ref{tab:ch6_test_y11}, \ref{tab:ch6_test_v8}, and~\ref{tab:ch6_test_deltas} present the locked 1,587-image official test results, followed by the unified 2$\times$2 visual synthesis in Figure~\ref{fig:ch6_ablation_val_test}.

\begin{table}[htbp]
\centering
\caption{GDUT-HWD Official Test Results (1,587 Images, 9,574 Boxes) for \texttt{YOLO11s} (\texttt{A0--A6}).}
\label{tab:ch6_test_y11}
\small
\renewcommand{\arraystretch}{1.18}
\resizebox{\textwidth}{!}{%
\begin{tabular}{|c|r|r|r|r|l|}
\hline
\rowcolor{slatehdr}
\textbf{Config} & \textbf{Precision (\%)} & \textbf{Recall (\%)} & \textbf{mAP@50 (\%)} & \textbf{mAP@50--95 (\%)} & \textbf{Tesla T4 Infer (ms)} \\
\hline
\textbf{A0 (S1)} & 91.0 & 82.8 & 88.8 & 58.5 & -- \\ \hline
\textbf{A1} & \textbf{91.2} & \textbf{84.4} & \textbf{89.8} & \textbf{59.1} & \textbf{7.5} \\ \hline
\textbf{A2} & 91.1 & 83.5 & 89.4 & 58.6 & -- \\ \hline
\textbf{A3} & 90.7 & 83.9 & 89.3 & 58.3 & -- \\ \hline
\textbf{A4} & 90.2 & 84.3 & 89.2 & 58.5 & -- \\ \hline
\textbf{A5} & 90.2 & 84.2 & 89.2 & 58.6 & -- \\ \hline
\rowcolor{slatehdr}
\textbf{A6} & 90.2 & 84.3 & 88.9 & 58.1 & 9.2 (unfused \texttt{RepConv}) \\ \hline
\end{tabular}}
\end{table}

\begin{table}[htbp]
\centering
\caption{GDUT-HWD Official Test Results (1,587 Images, 9,574 Boxes) for \texttt{YOLOv8s} (\texttt{A0--A6}).}
\label{tab:ch6_test_v8}
\small
\renewcommand{\arraystretch}{1.18}
\resizebox{\textwidth}{!}{%
\begin{tabular}{|c|r|r|r|r|l|}
\hline
\rowcolor{slatehdr}
\textbf{Config} & \textbf{Precision (\%)} & \textbf{Recall (\%)} & \textbf{mAP@50 (\%)} & \textbf{mAP@50--95 (\%)} & \textbf{Tesla T4 Infer (ms)} \\
\hline
\textbf{A0 (S1)} & 90.1 & 81.8 & 87.5 & 57.2 & -- \\ \hline
\textbf{A1} & 90.5 & 83.4 & \textbf{88.7} & \textbf{57.8} & \textbf{7.4} \\ \hline
\textbf{A2} & 90.6 & 83.1 & 88.6 & 57.2 & -- \\ \hline
\textbf{A3} & 89.7 & \textbf{83.6} & \textbf{88.7} & 57.2 & -- \\ \hline
\textbf{A4} & \textbf{90.9} & 82.6 & 88.3 & 57.2 & -- \\ \hline
\textbf{A5} & \textbf{90.9} & 82.4 & 88.3 & 57.3 & -- \\ \hline
\rowcolor{slatehdr}
\textbf{A6} & 90.3 & 82.6 & 88.4 & 57.4 & 9.1 (unfused \texttt{RepConv}) \\ \hline
\end{tabular}}
\end{table}

\begin{table}[htbp]
\centering
\caption{Official Test Deltas and Validation-to-Test Generalization Gap for \texttt{A6}.}
\label{tab:ch6_test_deltas}
\small
\renewcommand{\arraystretch}{1.16}
\resizebox{\textwidth}{!}{%
\begin{tabular}{|l|l|r|r|r|r|}
\hline
\rowcolor{slatehdr}
\textbf{Comparison} & \textbf{Backbone} & \textbf{$\Delta$ mAP@50} & \textbf{$\Delta$ mAP@50--95} & \textbf{Val$\to$Test mAP@50 Gap} & \textbf{Val$\to$Test mAP@50--95 Gap} \\
\hline
\texttt{A6 $-$ A0} & \texttt{YOLO11s} & $+0.1$ pp & $-0.4$ pp & $92.50\% \to 88.9\%$ ($-3.60$ pp) & $61.23\% \to 58.1\%$ ($-3.13$ pp) \\ \hline
\texttt{A6 $-$ A0} & \texttt{YOLOv8s} & $+0.9$ pp & $+0.2$ pp & $92.56\% \to 88.4\%$ ($-4.16$ pp) & $60.68\% \to 57.4\%$ ($-3.28$ pp) \\ \hline
\texttt{A6 $-$ A1} & \texttt{YOLO11s} & $-0.9$ pp & $-1.0$ pp & -- & -- \\ \hline
\texttt{A6 $-$ A1} & \texttt{YOLOv8s} & $-0.3$ pp & $-0.4$ pp & -- & -- \\ \hline
\end{tabular}}
\end{table}

\begin{figure}[htbp]
\centering
\includegraphics[width=0.98\textwidth]{figures/fig_gdut_ablation_val_test.png}
\caption{Complete \texttt{A0--A6} ablation comparison across Validation (316 images) and Official Test (1,587 images) splits on \texttt{YOLO11s} and \texttt{YOLOv8s}. Configuration \texttt{A1} consistently peaks in mAP@50 and mAP@50--95, while \texttt{A6} surpasses both \texttt{B1} and \texttt{B2} on test mAP@50.}
\label{fig:ch6_ablation_val_test}
\end{figure}

\noindent\textbf{Cross-Benchmark Standing Against SOTA Baselines \texttt{B1} and \texttt{B2}:} Comparing official 1,587-image test accuracy across our full model suite reveals a clear hierarchy: \texttt{YOLO11s A1} leads at \textbf{89.8\% mAP@50 / 59.1\% mAP@50--95}, followed by \texttt{YOLO11s A6} at \textbf{88.9\% / 58.1\%} and \texttt{YOLOv8s A1} at \textbf{88.7\% / 57.8\%}---all decisively outperforming reproduced \textbf{Baseline 1} (\texttt{YOLO11n-DCRCF}: $84.36\% / 54.37\%$, $+4.54$ to $+5.44$ pp mAP@50 gain) and surpassing \textbf{Baseline 2} (\texttt{YOLOv8s-DWR-ASPP-NWD}: $88.19\% / 57.98\%$) while using far fewer parameters ($9.49\text{M}$ deployed \texttt{YOLO11s A6} vs. $16.29\text{M}$ \texttt{B2}) and lower latency.

\subsubsection{Per-Class Test Analysis of A6 Across All Five Classes}

Table~\ref{tab:ch6_perclass_a6} and Figure~\ref{fig:ch6_perclass_a6} break down \texttt{A6}'s official 1,587-image test metrics by individual class. \texttt{yellow} helmets achieve the strongest AP@50 on both backbones ($92.7\%$ \texttt{YOLO11s}, $92.4\%$ \texttt{YOLOv8s}) due to high chromatic contrast against concrete/steel backgrounds. Conversely, \texttt{none} (bare heads) registers the lowest AP@50 ($84.6\%$ / $83.5\%$), AP@50--95 ($53.4\%$ / $51.8\%$), and Recall ($77.3\%$ / $75.1\%$) despite having the largest sample count in the test set (2,200 boxes) and 1,767 training boxes---ruling out simple class-frequency scarcity and pointing instead to extreme visual heterogeneity of uncovered heads against dark cluttered backgrounds. Notably, \texttt{YOLO11s A6} outperforms \texttt{YOLOv8s A6} most clearly on safety-critical \texttt{none} ($+1.1$ pp AP@50, $+1.6$ pp AP@50--95) and \texttt{red} ($+2.3$ pp AP@50), while \texttt{YOLOv8s A6} leads slightly on \texttt{white} ($90.4\%$ vs. $89.1\%$).

\begin{table}[htbp]
\centering
\caption{Per-Class Official Test Performance of \texttt{A6} Across All Five GDUT-HWD Classes (1,587 Images).}
\label{tab:ch6_perclass_a6}
\small
\renewcommand{\arraystretch}{1.18}
\resizebox{\textwidth}{!}{%
\begin{tabular}{|l|r|r|r|r|r|r|r|}
\hline
\rowcolor{slatehdr}
\textbf{Class} & \textbf{Test Boxes} & \multicolumn{3}{c|}{\textbf{YOLO11s A6 (\%)}} & \multicolumn{3}{c|}{\textbf{YOLOv8s A6 (\%)}} \\
\cline{3-8}
\rowcolor{slatehdr}
& & \textbf{Recall} & \textbf{AP@50} & \textbf{AP@50--95} & \textbf{Recall} & \textbf{AP@50} & \textbf{AP@50--95} \\
\hline
\texttt{blue}   & 1,361 & 86.6 & \textbf{89.8} & \textbf{59.4} & 83.3 & 89.3 & 58.8 \\ \hline
\texttt{white}  & 1,968 & 86.6 & 89.1 & 59.5 & \textbf{86.9} & \textbf{90.4} & \textbf{59.6} \\ \hline
\texttt{yellow} & 1,962 & \textbf{88.2} & \textbf{92.7} & \textbf{60.7} & 87.8 & 92.4 & 60.6 \\ \hline
\texttt{red}    & 2,083 & \textbf{82.7} & \textbf{88.6} & \textbf{57.4} & 80.1 & 86.3 & 56.0 \\ \hline
\texttt{none}   & 2,200 & \textbf{77.3} & \textbf{84.6} & \textbf{53.4} & 75.1 & 83.5 & 51.8 \\ \hline
\end{tabular}}
\end{table}

\begin{figure}[htbp]
\centering
\includegraphics[width=0.96\textwidth]{figures/fig_gdut_perclass_a6.png}
\caption{Per-class test AP@50 and AP@50--95 comparison between \texttt{YOLO11s A6} and \texttt{YOLOv8s A6} over 9,574 ground-truth bounding boxes, illustrating peak discrimination on \texttt{yellow} and the bare-head \texttt{none} bottleneck.}
\label{fig:ch6_perclass_a6}
\end{figure}

\subsubsection{Pre-Finetune Initialization Audit and Convergence Diagnostics}

Table~\ref{tab:ch6_init_epochs} and Figure~\ref{fig:ch6_init_epochs} present direct mechanistic evidence from \nolinkurl{initialization_check.json} (evaluated immediately after module patching and before Stage 2 gradient updates) alongside training convergence epochs from \texttt{results.csv}.

\begin{table}[htbp]
\centering
\caption{Pre-Finetune Function-Preservation Audit (\nolinkurl{initialization_check.json}) and Early-Stopping Epochs.}
\label{tab:ch6_init_epochs}
\resizebox{\textwidth}{!}{%
\renewcommand{\arraystretch}{1.18}
\begin{tabular}{|c|c|c|c|c|}
\hline
\rowcolor{slatehdr}
\textbf{Config} & \textbf{YOLO11s Pre-Finetune} & \textbf{YOLOv8s Pre-Finetune} & \textbf{YOLO11s Epochs} & \textbf{YOLOv8s Epochs} \\
\rowcolor{slatehdr}
& \textbf{mAP@50 / mAP@50--95 (\%)} & \textbf{mAP@50 / mAP@50--95 (\%)} & \textbf{(Stopped / Best)} & \textbf{(Stopped / Best)} \\
\hline
\textbf{A0} & 91.83 / 60.83 & 91.88 / 60.55 & Stage 1 & Stage 1 \\ \hline
\textbf{A1} & -- (stock init) & -- (stock init) & 70 / 20 & 63 / 13 \\ \hline
\textbf{A2} & 91.83 / 60.83 & 91.88 / 60.55 & 52 / 2  & 65 / 15 \\ \hline
\textbf{A3} & 91.83 / 60.83 & 91.88 / 60.55 & 66 / 16 & 51 / 1  \\ \hline
\textbf{A4} & -- (loss only) & -- (loss only) & 65 / 15 & 66 / 16 \\ \hline
\textbf{A5} & \textbf{86.34 / 56.15} & 91.88 / 60.55 & 65 / 15 & 51 / 1  \\ \hline
\textbf{A6} & \textbf{86.34 / 56.15} & 91.88 / 60.55 & 66 / 16 & 66 / 16 \\ \hline
\end{tabular}
}
\end{table}

\begin{figure}[htbp]
\centering
\includegraphics[width=0.98\textwidth]{figures/fig_gdut_init_check_and_epochs.png}
\caption{Mechanistic convergence diagnostics: (a) Exact function preservation for \texttt{A2}/\texttt{A3} on both backbones and \texttt{A5}/\texttt{A6} on \texttt{YOLOv8s}, versus the $-5.49$ pp pre-finetune drop on \texttt{YOLO11s} caused by replacing \texttt{C2PSA}; (b) Real validation mAP@50--95 convergence trajectories from \texttt{results.csv} confirming clean early stopping (\texttt{patience = 50}).}
\label{fig:ch6_init_epochs}
\end{figure}

Crucially, Figure~\ref{fig:ch6_init_epochs}(b) refutes any claim that \texttt{A6} underperformed \texttt{A1} merely due to premature stopping: both \texttt{YOLO11s A6} and \texttt{YOLOv8s A6} peaked at epoch 16, trained 50 additional epochs to epoch 66 under \texttt{patience = 50}, and gradually drifted downward in validation mAP@50--95 ($61.23\% \to 59.96\%$ on \texttt{YOLO11s}; $60.69\% \to 59.31\%$ on \texttt{YOLOv8s}), confirming that the saved epoch-16 \texttt{best.pt} represents the true optimum under this recipe. Figures~\ref{fig:ch6_a6_curves}, \ref{fig:ch6_test_curves}, and~\ref{fig:ch6_qual_preds} present authentic Kaggle training curves, PR/confusion matrices, and qualitative prediction grids directly from \texttt{nvthanh2004}'s canonical runs.

\begin{figure}[htbp]
\centering
\includegraphics[width=0.96\textwidth]{figures/a6_yolo11s_results.png}
\caption{Authentic Kaggle Stage 2 training and validation loss/metric curves (\texttt{results.png}) for \texttt{GDUT A6 YOLO11s} (\texttt{ce2a883e...}), showing rapid recovery from the epoch-0 \texttt{C2PSA} drop and stable early stopping at epoch 66.}
\label{fig:ch6_a6_curves}
\end{figure}

\begin{figure}[htbp]
\centering
\begin{subfigure}[b]{0.48\textwidth}
  \centering
  \includegraphics[width=\textwidth]{figures/a6_yolo11s_val_pr.png}
  \caption{A6 \texttt{YOLO11s} Validation PR Curve (mAP@50 92.5\%)}
\end{subfigure}
\hfill
\begin{subfigure}[b]{0.48\textwidth}
  \centering
  \includegraphics[width=\textwidth]{figures/a6_yolo11s_val_cm.png}
  \caption{A6 \texttt{YOLO11s} Validation Confusion Matrix}
\end{subfigure}\\[6pt]
\begin{subfigure}[b]{0.48\textwidth}
  \centering
  \includegraphics[width=\textwidth]{figures/a1_yolo11s_test_pr.png}
  \caption{A1 \texttt{YOLO11s} Official Test PR Curve (mAP@50 89.8\%)}
\end{subfigure}
\hfill
\begin{subfigure}[b]{0.48\textwidth}
  \centering
  \includegraphics[width=\textwidth]{figures/a1_yolo11s_test_cm.png}
  \caption{A1 \texttt{YOLO11s} Official Test Confusion Matrix}
\end{subfigure}
\caption{Authentic Kaggle PR curves and normalized confusion matrices comparing \texttt{A6 YOLO11s} on Validation (top) and top-performing \texttt{A1 YOLO11s} on the 1,587-image Official Test split (bottom). Across both splits, inter-color helmet confusion remains minimal ($<2\%$), while the dominant error mode is \texttt{none} $\leftrightarrow$ background leakage.}
\label{fig:ch6_test_curves}
\end{figure}

\begin{figure}[htbp]
\centering
\begin{subfigure}[b]{0.49\textwidth}
  \centering
  \includegraphics[width=\textwidth]{figures/a6_yolo11s_val_pred0.jpg}
  \caption{A6 \texttt{YOLO11s} Validation Predictions (\texttt{val\_batch0})}
\end{subfigure}
\hfill
\begin{subfigure}[b]{0.49\textwidth}
  \centering
  \includegraphics[width=\textwidth]{figures/a1_yolo11s_test_pred0.jpg}
  \caption{A1 \texttt{YOLO11s} Official Test Predictions (\texttt{test\_batch0})}
\end{subfigure}
\caption{Authentic qualitative five-class helmet and bare-head detections from canonical GDUT Kaggle evaluations, demonstrating crisp multi-color discrimination across crowded construction sites and dense scaffolding.}
\label{fig:ch6_qual_preds}
\end{figure}

\subsection{Inference Optimization (Before vs. After) and Edge Hardware Profiling}

\subsubsection{Progressive RTX 3050 Optimization Ladder (GDUT-A6)}

To fulfill Review 3 Stage 1 Items 1.4 and 1.5 with verifiable hardware measurements, we benchmarked our deployed \texttt{Rep-YOLO11s-GDUT-A6} checkpoint across a five-stage optimization ladder on an edge-class \textbf{NVIDIA GeForce RTX 3050 Laptop GPU} (4 GB VRAM, PyTorch \texttt{2.0.0+cu118}) in \nolinkurl{Output/task3_eval_opt/optimization_benchmark_table.json}. Table~\ref{tab:ch6_opt_ladder} demonstrates that replacing Python region loops in \texttt{BiFormerBlockLite} with vectorized \texttt{torch.gather}, enabling CUDA AMP FP16 over FP32 master weights, and applying micro-batching ($N=4$) reduces per-frame inference latency from \textbf{12.86 ms (77.8 FPS)} to \textbf{6.81 ms/frame (146.9 FPS, $1.89\times$ GPU speedup; $46.4\times$ over 4-thread CPU)} with \textbf{0 box-count deviation} verified across 12 real construction video frames (\nolinkurl{fp32_boxes = 2}, \nolinkurl{fp16_amp_boxes = 2}, \nolinkurl{max_abs_count_delta = 0}).

\begin{table}[htbp]
\centering
\caption{Before vs. After Inference Optimization Benchmark on \texttt{Rep-YOLO11s-GDUT-A6} (RTX 3050 Laptop GPU, $640 \times 640$).}
\label{tab:ch6_opt_ladder}
\resizebox{\textwidth}{!}{%
\renewcommand{\arraystretch}{1.20}
\begin{tabular}{|l|l|c|r|r|r|r|}
\hline
\rowcolor{slatehdr}
\textbf{Optimization Stage} & \textbf{Hardware / Precision} & \textbf{Batch} & \textbf{Latency (ms/f)} & \textbf{FPS} & \textbf{VRAM (MB)} & \textbf{Speedup vs. Before} \\
\hline
CPU Reference (4-Thread) & CPU FP32 & 1 & 315.69 & 3.2 & 0.0 & $0.04\times$ \\ \hline
\textbf{Before} (Loop BiFormer) & RTX 3050 FP32 & 1 & 12.86 & 77.8 & 282.3 & $1.00\times$ \\ \hline
\textbf{After-1} (+ Vec \texttt{torch.gather} + RepConv) & RTX 3050 FP32 & 1 & 11.93 & 83.8 & 282.3 & $1.08\times$ \\ \hline
\textbf{After-2} (+ CUDA AMP FP16) & RTX 3050 FP16-AMP & 1 & 10.50 & 95.2 & \textbf{278.0} & $1.22\times$ \\ \hline
\rowcolor{slatehdr}
\textbf{After-3 (Micro-Batch $N=4$)} & \textbf{RTX 3050 FP16-AMP} & \textbf{4} & \textbf{6.81} & \textbf{146.9} & 341.0 & \textbf{$1.89\times$} \\ \hline
\end{tabular}
}
\end{table}

\subsubsection{Dedicated Verification of \nolinkurl{RepConv.switch_to_deploy()} Equivalence and Speedup}

As noted transparently in Section 9 of \nolinkurl{01_KET_QUA_A1_A6.md}, the Kaggle \nolinkurl{GDUT_final_test_evaluation.ipynb} log reported \texttt{A6} at \texttt{9.2 ms} versus \texttt{A1} at \texttt{7.5 ms} on Tesla T4 \textit{without calling \nolinkurl{switch_to_deploy()}}. To close this review gap quantitatively with real hardware proof, we executed a dedicated before-vs-after verification script on the local unfused \nolinkurl{GDUT_A6_YOLO11s_best.pt} checkpoint (identical graph with six stride-2 \texttt{RepConv} layers at \texttt{[1, 3, 5, 7, 17, 20]}; \nolinkurl{allow_tf32 = False}, 15 warmup + 60 timed iterations on RTX 3050, saved to \nolinkurl{Output/task3_eval_opt/gdut_a6_switch_to_deploy_benchmark.json}). Table~\ref{tab:ch6_repconv_deploy} confirms exact parameter reduction (\textbf{$-300,992$ parameters}, $-3.07\%$), tight numerical equivalence across the full $1 \times 9 \times 8400$ decoded output tensor (\textbf{mean absolute deviation $6.28 \times 10^{-6}$}, max absolute deviation $7.93 \times 10^{-4}$ over 640-pixel coordinates), and real wall-clock latency reductions of \textbf{$-3.6\%$ in FP32} ($24.870 \to 23.973$ ms) and \textbf{$-6.3\%$ in CUDA AMP FP16} ($11.521 \to 10.797$ ms, lifting throughput from $86.79$ to \textbf{92.62 FPS}).

\begin{table}[H]
\centering
\caption{Dedicated RTX 3050 Verification of \nolinkurl{RepConv.switch_to_deploy()} Across All Six Stride-2 Blocks (\texttt{layers 1, 3, 5, 7, 17, 20}) on Unfused \nolinkurl{GDUT_A6_YOLO11s_best.pt}.}
\label{tab:ch6_repconv_deploy}
\resizebox{\textwidth}{!}{%
\renewcommand{\arraystretch}{1.20}
\begin{tabular}{|l|r|r|r|}
\hline
\rowcolor{slatehdr}
\textbf{Measured Metric (B=1, $640\times 640$, 60 Iters)} & \textbf{Before \nolinkurl{switch_to_deploy()}} & \textbf{After \nolinkurl{switch_to_deploy()}} & \textbf{Verified Delta} \\
\hline
Unfused / Fused \texttt{RepConv} Count & 6 multi-branch blocks & 6 single $3\times 3$ blocks & 100\% fused \\ \hline
Total Model Parameters & 9,788,575 (9.789M) & \textbf{9,487,583 (9.488M)} & \textbf{$-300,992$ ($-3.07\%$)} \\ \hline
IEEE FP32 Output Parity (\texttt{mean / max abs diff}) & Reference & $6.28 \times 10^{-6}$ / $7.93 \times 10^{-4}$ & Mathematically equivalent \\ \hline
FP32 Latency / Throughput & 24.870 ms / 40.21 FPS & \textbf{23.973 ms / 41.71 FPS} & $-0.897$ ms ($-3.6\%$) \\ \hline
CUDA AMP FP16 Latency / Throughput & 11.521 ms / 86.79 FPS & \textbf{10.797 ms / 92.62 FPS} & \textbf{$-0.724$ ms ($-6.3\%$, $+5.83$ FPS)} \\ \hline
Peak Allocated VRAM (FP32 / AMP FP16) & 167.44 / 216.04 MB & 169.99 / 222.35 MB & Negligible ($<7$ MB) \\ \hline
\end{tabular}
}
\end{table}

\noindent\textbf{Technical Reconciliation of FP32 Latency Between Table~\ref{tab:ch6_opt_ladder} (\texttt{11.93 ms}) and Table~\ref{tab:ch6_repconv_deploy} (\texttt{23.973 ms}):} Notice that single-batch FP32 latency in Table~\ref{tab:ch6_repconv_deploy} ($23.973\text{ ms}$) is approximately double the \texttt{After-1} FP32 figure in Table~\ref{tab:ch6_opt_ladder} ($11.93\text{ ms}$). This difference stems strictly from CUDA precision-backend flags and runtime serialization rather than model discrepancy: (1) Table~\ref{tab:ch6_opt_ladder} benchmarked the exported \nolinkurl{rep_yolo11s_gdut_a6_deploy.pt} checkpoint under PyTorch \texttt{2.0.0}'s default configuration (\nolinkurl{cudnn.allow_tf32 = True}), which routes cuDNN $3\times 3$ convolutions onto Ampere Tensor Cores using 19-bit TensorFloat-32 (TF32) arithmetic alongside vectorized \texttt{torch.gather} in \nolinkurl{custom_ablation_modules.BiFormerBlockLite}; (2) conversely, Table~\ref{tab:ch6_repconv_deploy} explicitly enforced \nolinkurl{allow_tf32 = False} for both cuDNN and cuBLAS on the raw \nolinkurl{GDUT_A6_YOLO11s_best.pt} training checkpoint so that output parity ($6.28 \times 10^{-6}$ mean abs diff) reflects strict IEEE-754 23-bit mantissa FP32 execution on standard SIMT cores without TF32 mantissa truncation. Re-benchmarking \nolinkurl{GDUT_A6_YOLO11s_best.pt} on the identical RTX 3050 under \nolinkurl{cudnn.allow_tf32 = True} confirms a \texttt{switch\_to\_deploy()} speedup from \textbf{15.145 ms (66.03 FPS)} to \textbf{14.107 ms (70.89 FPS, $-6.85\%$ latency)}, reconciling both benchmarks transparently (with the minor $2.18\text{ ms}$ residual vs. $11.93\text{ ms}$ attributable to laptop GPU thermal/clock variance and raw training wrapper vs. standalone exported \texttt{DetectionModel}).

\subsubsection{Multi-Stream 1080p Edge Scaling and Embedded Roofline Projection}

Table~\ref{tab:ch6_multistream} reports full end-to-end batched multi-stream profiling (1080p decode + GPU inference + post-processing) of the deployed \texttt{Rep-YOLO11s-GDUT-A6} checkpoint (\nolinkurl{rep_yolo11s_gdut_a6_deploy.pt}) on our RTX 3050 edge node (\nolinkurl{Output/operational_multistream/rtx3050_gdut_a6_batched_multistream_results.json}) alongside analytical Roofline guidance for embedded Jetson-class deployments.

\begin{table}[htbp]
\centering
\caption{End-to-End 1080p Batched Multi-Stream Profiling of \texttt{Rep-YOLO11s-GDUT-A6} on RTX 3050 Edge Node ($N=1,2,4,8$).}
\label{tab:ch6_multistream}
\resizebox{\textwidth}{!}{%
\renewcommand{\arraystretch}{1.18}
\begin{tabular}{|c|r|r|r|r|r|r|r|c|}
\hline
\rowcolor{slatehdr}
\textbf{Streams ($N$)} & \textbf{Agg. FPS} & \textbf{FPS / Stream} & \textbf{Mean (ms)} & \textbf{p95 (ms)} & \textbf{GPU Util (\%)} & \textbf{VRAM (MB)} & \textbf{Power (W)} & \textbf{25 FPS SLA} \\
\hline
1 & 40.9 & \textbf{40.9} & 23.98 & 28.61 & 33.3 & 446.0 & 31.1 & \textbf{PASS} \\ \hline
2 & 51.5 & \textbf{25.8} & 37.71 & 42.32 & 35.0 & 450.0 & 38.2 & \textbf{PASS} \\ \hline
4 & 55.1 & 13.8 & 70.34 & 77.20 & 37.1 & 578.0 & 37.7 & Below 25 FPS (10--15 FPS CCTV ok) \\ \hline
8 & \textbf{57.1} & 7.1  & 135.11 & 146.87 & 42.8 & 706.0 & 38.9 & Requires stride-2 / multi-node \\ \hline
\end{tabular}
}
\end{table}

Memory consumption stays remarkably compact---under \textbf{706 MB VRAM} even at 8 simultaneous 1080p streams (with $N=1$ at 40.9 FPS and $N=2$ at 25.8 FPS/stream meeting the 25 FPS SLA)---meaning an NVIDIA Jetson Orin Nano (8 GB unified memory, 40 TOPS INT8 / 20 TFLOPS FP16, 15W TDP) or Xavier NX easily accommodates the full single-container Docker stack with ample headroom for 1--2 full-rate 25 FPS channels or 4 channels at 12--15 FPS industrial CCTV sampling.

\subsection{Software Verification, API Load/Stress Testing, Field Audit, and UAT/SUS}

\begin{enumerate}[leftmargin=1.8em, itemsep=4pt]
  \item \textbf{Automated Unit and Integration Verification (\texttt{pytest}):} Executing \texttt{python -m pytest tests -q} on CPU (\nolinkurl{CUDA_VISIBLE_DEVICES=-1}) passes \textbf{108 / 108 tests (100\%)} in 24.98 seconds (and passes 108/108 in weightless CI stub mode \nolinkurl{MODEL_PATH=ci_stub_missing.pt}). Test coverage spans Dynamic Taxonomy GDUT adaptation (\texttt{cls=2} $\to$ \texttt{yellow}), 100-frame stationary-worker PPE retention, ByteTrack/Re-ID anti-hijack, 8-frame violation debouncing, FastAPI endpoints (\texttt{/api/metrics}, \texttt{/metrics}, \nolinkurl{/api/predict_image}, \texttt{/api/feedback}, \nolinkurl{/api/switch_model}), and mocked Telegram/Zalo dispatchers.
  \item \textbf{FastAPI Concurrency Load and Stress Testing:} Benchmarking the live ASGI server via \nolinkurl{scripts/load_test_api.py} (\nolinkurl{Output/load_test_results.json}) across $C \in \{10, 25, 50, 100\}$ concurrent users (1,400 total HTTP requests across telemetry, status, alerts, and workers endpoints) demonstrates sustained throughput peaking at \textbf{346.0 RPS} ($C=25$) and holding \textbf{320.4 RPS} at $C=100$ with \textbf{0 HTTP errors (0.0\% failure rate)} across every concurrency tier (Table~\ref{tab:ch6_load_test} and Figure~\ref{fig:ch6_opt_load}).
  \item \textbf{Real-World Construction CCTV Video Testing:} Verified smooth end-to-end execution across our 11-video repository in \texttt{video\_test/} spanning 1080p and 4K UHD construction scenes (\nolinkurl{construction_site_workers_1080p.mp4}, \nolinkurl{10810476-hd_1920_1080_30fps.mp4}, \nolinkurl{13921040-uhd_3840_2160_30fps.mp4}, \texttt{brickmasonry 11.mp4}, \texttt{concreting 10.mp4}), confirming stable real-time HUD telemetry, instant hot-swap via \nolinkurl{switch_to_gdut.bat}, and zero crashes under 4K automatic downscaling.
  \item \textbf{User Acceptance Testing (UAT) and SUS Protocol:} A standardized 10-item System Usability Scale (SUS) questionnaire~\cite{brooke1996sus} (5-point Likert, $\text{SUS} = 2.5 \sum c_i \in [0, 100]$, target $\ge 68.0$) paired with four timed operational scenarios (Camera Switch; GDUT 5-Class Hot-Swap via \nolinkurl{/api/switch_model} verifying \nolinkurl{is_helmet_only=true}, String-First 4-color helmet distribution, and \texttt{none} violation alerting [with \textit{Solo Focus} activating only when hot-swapped to a \texttt{person}-class model]; Evidence Lightbox + HITL Feedback; and Image Audit Upload) is packaged in \nolinkurl{docs/UAT_SURVEY_TEMPLATE_AND_SUS_PROTOCOL.md} and reproduced verbatim in Appendix D. Consistent with our zero-fabrication commitment, because external construction-officer field trials ($N \ge 10$) are scheduled post-defense, no synthetic SUS score is fabricated here.
\end{enumerate}

\begin{table}[htbp]
\centering
\caption{FastAPI Concurrency Load and Stress Benchmark (\nolinkurl{scripts/load_test_api.py}, \nolinkurl{Output/load_test_results.json}).}
\label{tab:ch6_load_test}
\small
\renewcommand{\arraystretch}{1.18}
\resizebox{\textwidth}{!}{%
\begin{tabular}{|c|r|r|r|r|r|r|r|r|}
\hline
\rowcolor{slatehdr}
\textbf{Concurrency ($C$)} & \textbf{Requests} & \textbf{Wall (s)} & \textbf{Throughput (RPS)} & \textbf{Avg (ms)} & \textbf{p50 (ms)} & \textbf{p95 (ms)} & \textbf{p99 (ms)} & \textbf{Error (\%)} \\
\hline
10  & 200 & 0.737 & 271.3 & 35.24  & 27.59  & 69.26  & 80.25  & \textbf{0.0} \\ \hline
25  & 300 & 0.867 & \textbf{346.0} & 68.02  & 52.57  & 125.08 & 129.32 & \textbf{0.0} \\ \hline
50  & 400 & 1.265 & 316.3 & 147.96 & 137.76 & 257.07 & 294.87 & \textbf{0.0} \\ \hline
100 & 500 & 1.561 & 320.4 & 277.46 & 266.44 & 371.11 & 436.89 & \textbf{0.0} \\ \hline
\end{tabular}}
\end{table}

\begin{figure}[htbp]
\centering
\includegraphics[width=0.98\textwidth]{figures/fig_gdut_optimization_and_multistream.png}
\caption{Production optimization and load verification: (a) Progressive RTX 3050 speedup of \texttt{Rep-YOLO11s-GDUT-A6} reaching 146.9 FPS ($1.89\times$, $5.9\times$ above 25 FPS SLA); (b) FastAPI concurrent load throughput peaking at 346.0 RPS with 0.0\% errors up to $C=100$.}
\label{fig:ch6_opt_load}
\end{figure}

\section{Discussion}

Synthesizing all quantitative evidence allows us to answer \textbf{RQ1--RQ4} candidly and equip the team for deep technical defense before the Review 3 committee:

\begin{itemize}[leftmargin=1.8em, itemsep=4pt]
  \item \textbf{Answer to RQ1 (Deduplication and Color-Safe Augmentation):} Removing the six exact SHA-256 pixel duplicates between \texttt{trainval} and \texttt{test} eliminated silent data leakage while keeping all 1,587 official test images intact. Pairing offline color-preserving HardAug (\nolinkurl{hsv_h=0, hsv_s=0}) with online \texttt{mixup=0.1} (\texttt{A1}) delivered the largest single accuracy leap in the entire study: $+1.0$ pp test mAP@50 ($88.8\% \to 89.8\%$) and $+0.6$ pp test mAP@50--95 ($58.5\% \to 59.1\%$) on \texttt{YOLO11s}, and $+1.2$ pp / $+0.6$ pp on \texttt{YOLOv8s}, while holding inter-color confusion below $2\%$.
  \item \textbf{Answer to RQ2 (Fair SOTA Baseline Audit):} Neither published baseline outperformed its matched stock YOLO counterpart on five-class GDUT-HWD. \texttt{B1} (\texttt{YOLO11n-DCRCF}) lost $-1.529$ pp to Stock \texttt{YOLO11n} across 5 seeds because 73.78\% randomly initialized neck parameters destabilized P3 confidence calibration (causing 84/98 \texttt{none} misses to fall below $0.25$). \texttt{B2} (\texttt{YOLOv8s-DWR-ASPP-NWD}) lost $-0.633$ pp to Stock \texttt{YOLOv8s} across 5 seeds while inflating FLOPs by $+45.0\%$. These findings demonstrate why reproducing published SOTA architectures under matched multi-seed controls is indispensable.
  \item \textbf{Answer to RQ3 (Why \texttt{A1} Outperformed \texttt{A6} on GDUT-HWD):} Four empirical factors explain why combining all four custom components in \texttt{A6} did not yield strictly additive accuracy gains over \texttt{A1}:
  \begin{enumerate}[leftmargin=1.6em, itemsep=2pt]
    \item \textit{Online Augmentation Asymmetry:} \texttt{A1} benefited from \texttt{mixup = 0.1} and \texttt{hsv\_v = 0.45}, whereas \texttt{A2--A6} used \texttt{mixup = 0.0} and \texttt{hsv\_v = 0.40}, confounding pure architectural attribution.
    \item \textit{Pre-Finetune \texttt{C2PSA} Disruption on \texttt{YOLO11s}:} \nolinkurl{initialization_check.json} proves directly that replacing pretrained \texttt{C2PSA} with identity-initialized \texttt{BiFormerBlockLite} immediately dropped \texttt{YOLO11s} validation mAP@50 from $91.83\%$ to $86.34\%$ ($-5.49$ pp). Fine-tuning had to spend capacity recovering lost \texttt{C2PSA} representations before improving further. In sharp contrast, on \texttt{YOLOv8s}---where \texttt{SPPF} remained intact and \texttt{BiFormer} was appended after it---pre-finetune accuracy held steady at $91.88\%$, allowing \texttt{YOLOv8s A6} to achieve the highest validation precision across all 14 models (\textbf{93.21\%}) and net gains over \texttt{A0} on both Val ($+0.68$ pp) and Test ($+0.9$ pp).
    \item \textit{Focal-EIoU Localization Trade-off:} On tiny ambiguous bare heads (\texttt{none}), down-weighting low-IoU regressions via $\mathrm{IoU}^{0.5}$ did not surpass CIoU in mAP@50--95 (\texttt{A4} matched \texttt{A0} on test mAP@50--95 at $58.5\%$ / $57.2\%$).
    \item \textit{Role of Scientific Ablation:} An ablation study exists to test design hypotheses objectively---not to force a maximal-complexity model to win every column. On a 1,265-scene dataset, clean data engineering (\texttt{A1}) often dominates added architectural parameters.
  \end{enumerate}
  \item \textbf{Answer to RQ4 (Edge Optimization and VMS Reliability):} Vectorizing \texttt{BiFormerBlockLite}, verifying \nolinkurl{RepConv.switch_to_deploy()} ($-300,992$ params, $-6.3\%$ AMP FP16 latency, $6.28 \times 10^{-6}$ mean diff), enabling CUDA AMP FP16, and micro-batching ($N=4$) drove \texttt{Rep-YOLO11s-GDUT-A6} to \textbf{146.9 FPS} ($6.81$ ms/frame) on RTX 3050 with zero detection drift, while AEGIS-CCTV passed 108/108 tests and sustained 346.0 RPS with 0.0\% HTTP failures.
\end{itemize}

\section{Recommendations}

Based directly on verified empirical evidence rather than theoretical preference, we offer four actionable engineering recommendations for construction-site deployment and follow-on research:
\begin{enumerate}[leftmargin=1.8em, itemsep=3pt]
  \item \textbf{Production Checkpoint Selection on GDUT-HWD:} Deploy \textbf{\texttt{YOLO11s A1}} when maximizing overall detection and localization accuracy ($89.8\%$ test mAP@50, $59.1\%$ test mAP@50--95, $7.5$ ms on T4). Deploy \textbf{\texttt{YOLOv8s A6}} in low-false-alarm environments prioritizing peak precision ($93.21\%$ validation precision), or \textbf{\texttt{YOLO11s A6}} (after calling \nolinkurl{switch_to_deploy()}) when demonstrating the full multi-module re-parameterized research architecture.
  \item \textbf{Mandatory \nolinkurl{switch_to_deploy()} Before Edge Export:} Always execute \nolinkurl{RepConv.switch_to_deploy()} followed by strict FP32 output parity verification before exporting ONNX/TensorRT engines so all six stride-2 training branches collapse cleanly into single $3\times 3$ kernels.
  \item \textbf{Targeted Remedies for Bare-Head (\texttt{none}) Recall:} Because \texttt{none} remains the weakest class ($84.6\%$ AP@50) due to background confusion and $<32$ px scale rather than sample scarcity, future datasets should prioritize hard-negative mining (caps, hoods, shadows) and higher-resolution tiling ($960\text{--}1280$ px) on distant camera zones.
  \item \textbf{Factorial Multi-Seed Protocol for Next-Stage Research:} Future architectural studies on GDUT-HWD should lock identical online augmentation across \texttt{A1--A6}, append \texttt{BiFormer} after \texttt{C2PSA} rather than replacing it on \texttt{YOLO11s}, and execute $\ge 3$--$5$ random seeds per variant to quantify confidence intervals.
\end{enumerate}

\section{Personal Reflection}

Executing Capstone Project \textbf{FA26AI16} across Reviews 1, 2, and 3 transformed our understanding of applied Artificial Intelligence engineering. Early in the project, we assumed that adding sophisticated published blocks (\texttt{CoordConv}, \texttt{RepConv}, \texttt{BiFormer}, \texttt{Focal-EIoU}) or merging larger datasets would automatically guarantee higher accuracy. Through rigorous forensic auditing---catching six SHA-256 pixel duplicates in \texttt{GDUT-HWD}, witnessing \texttt{C2PSA} replacement drop pre-finetune accuracy by $-5.49$ pp, proving that two published SOTA baselines underperformed stock YOLO models across five seeds, and debugging subtle production faults like the GDUT \texttt{cls=2} yellow-to-blue color bug and stationary-worker suppression---we learned that \textbf{data integrity, controlled instrumentation, honest negative analysis, and resilient systems engineering matter far more than architectural complexity alone}. These lessons directly prepared Group GFA26AI17 for real-world industrial AI engineering.
